\documentclass[11pt,a4paper]{article}
\usepackage[margin=2cm]{geometry}
\usepackage{amsmath,amssymb}
\usepackage{xcolor}
\usepackage{enumitem}
\usepackage{tcolorbox}
\usepackage{hyperref}
\usepackage{array}
\usepackage{multicol}
\tcbuselibrary{breakable,skins}

\definecolor{hi}{HTML}{D63A6F}
\definecolor{accent}{HTML}{1F4E79}
\definecolor{warn}{HTML}{B91C1C}
\hypersetup{colorlinks=true,linkcolor=accent,urlcolor=accent}

\newtcolorbox{must}{colback=hi!8,colframe=hi!75!black,boxrule=0.7pt,arc=2pt,
  fonttitle=\bfseries,title=Algorithm,breakable}
\newtcolorbox{useful}{colback=accent!6,colframe=accent!60!black,boxrule=0.6pt,arc=2pt,
  fonttitle=\bfseries,title=Worth knowing,breakable}
\newtcolorbox{gotcha}{colback=warn!6,colframe=warn!70!black,boxrule=0.6pt,arc=2pt,
  fonttitle=\bfseries,title=Exam trap,breakable}
\newtcolorbox{plain}[1]{colback=gray!5,colframe=gray!50!black,boxrule=0.5pt,arc=2pt,
  fonttitle=\bfseries,title=#1,breakable}

\newcommand{\seq}{\Rightarrow}
\newcommand{\imp}{\rightarrow}
\newcommand{\eq}{\simeq}
\newcommand{\T}{\mathbf{t}}
\newcommand{\F}{\mathbf{f}}
\newcommand{\rn}[1]{\;({#1})}

\setlength{\parindent}{0pt}
\setlength{\parskip}{4pt}

\title{\vspace{-1.5cm}\textbf{Logic and Proof} \\[2pt] \large Algorithms How-To + One-Page Rule Reference}
\author{\small Part IB --- Cambridge CST --- Paper 6}
\date{}

\begin{document}
\maketitle
\vspace{-0.8cm}

\section*{What the exam actually asks}

Logic \& Proof is an \textbf{algorithm-running} course: questions hand you a formula or clause set and say ``apply DPLL / resolution / build the BDD.'' Across \textbf{2018--2025 (16 questions)}:

\begin{center}
\begin{tabular}{lcc}
\hline
\textbf{Topic} & \textbf{Appearances} & \textbf{Share} \\
\hline
Resolution theorem proving & 7 & 44\% \\
Binary Decision Diagrams & 6 & 38\% \\
DPLL method & 5 & 31\% \\
First-order logic / Skolemisation \& unification & 3 & 19\% \\
Sequent calculus & 3 & 19\% \\
CNF conversion; Fourier--Motzkin; model checking; S4/S5 modal & 2 each & 12\% \\
\hline
\end{tabular}
\end{center}

\textbf{80/20 verdict.} \textbf{Resolution, BDDs and DPLL are the spine ($\sim$60\% of marks)} --- you must be able to \emph{run each by hand, showing every step}. They all start from \textbf{clause/CNF form}, so master that conversion first. Then sequent calculus + quantifier rules, unification, Fourier--Motzkin, and modal S4. This guide gives the \textbf{recipe + a worked example} for each algorithm, then a \textbf{one-page rule reference} (last page) to memorise.

\hrulefill

\section{Clause / CNF conversion (everything starts here)}

\begin{must}[title=Propositional CNF (4 steps)]
\begin{enumerate}[leftmargin=*,itemsep=0pt]
  \item \textbf{Eliminate $\leftrightarrow,\imp$:}\ \ $A\leftrightarrow B \eq (A\imp B)\wedge(B\imp A)$, \ \ $A\imp B\eq \neg A\vee B$.
  \item \textbf{Push $\neg$ inward} (to atoms) $\Rightarrow$ \emph{NNF}: \ $\neg\neg A\eq A$,\ \ de Morgan $\neg(A\wedge B)\eq\neg A\vee\neg B$, etc.
  \item \textbf{Push $\vee$ inward} (distribute over $\wedge$): \ $A\vee(B\wedge C)\eq(A\vee B)\wedge(A\vee C)$.
  \item \textbf{Simplify:} drop a disjunction with $P$ and $\neg P$ ($\eq\T$); $(P\vee A)\wedge(\neg P\vee A)\eq A$; remove repeats.
\end{enumerate}
A \textbf{clause} is a disjunction of literals, written as a set $\{\neg K_1,\dots,L_1,\dots\}$; the empty clause $\square=\{\,\}$ is \textbf{false} (contradiction).
\end{must}

\begin{must}[title=First-order clause form (for resolution)]
\begin{enumerate}[leftmargin=*,itemsep=0pt]
  \item \textbf{Negate} the goal (resolution and DPLL are \emph{refutation} methods --- prove unsatisfiability).
  \item Eliminate $\leftrightarrow,\imp$; push $\neg$ in to get \textbf{NNF}.
  \item \textbf{Skolemise:} for each $\exists y$ inside $\forall x_1\dots\forall x_k$, replace $y$ by $f(x_1,\dots,x_k)$ for a \emph{new} $f$ (a Skolem \emph{constant} if no enclosing $\forall$). Push quantifiers \emph{in} first (miniscope) $\Rightarrow$ fewer/smaller Skolem args $\Rightarrow$ shorter proofs.
  \item \textbf{Drop the $\forall$'s} (clauses are implicitly universally quantified).
  \item Distribute to \textbf{CNF}; write each conjunct as a clause (set of literals).
\end{enumerate}
Skolemisation preserves \emph{(un)satisfiability}, not validity --- which is exactly what refutation needs.
\end{must}

\section{DPLL (propositional SAT) --- 31\%}

\begin{must}
Refute a clause set (show unsatisfiable). Repeat:
\begin{enumerate}[leftmargin=*,itemsep=0pt]
  \item \textbf{Delete tautologies} --- any clause containing $P$ and $\neg P$.
  \item \textbf{Unit propagation} --- for each unit clause $\{L\}$: delete every clause containing $L$, and delete $\neg L$ from all clauses.
  \item \textbf{Pure-literal deletion} --- if $L$ occurs but $\neg L$ never does, delete all clauses containing $L$.
  \item \textbf{Stop:} empty clause $\square$ generated $\Rightarrow$ \textbf{unsatisfiable} (refutation found). All clauses deleted $\Rightarrow$ \textbf{satisfiable}.
  \item \textbf{Case split} on a literal $L$: recurse on the set $+\{L\}$ and on the set $+\{\neg L\}$. Satisfiable iff \emph{either} branch is.
\end{enumerate}
\end{must}

\begin{plain}{Worked: is $P\vee Q\imp Q\vee R$ a theorem?}
Negate $+$ CNF gives clauses $\{P,Q\}$, $\{\neg Q\}$, $\{\neg R\}$.
\[
\begin{array}{llll}
\{P,Q\} & \{\neg Q\} & \{\neg R\} & \text{initial}\\
\{P\}   &           & \{\neg R\} & \text{unit }\neg Q\\
        &           &            & \text{units }P,\ \neg R\ \text{(also pure)} \Rightarrow\ \text{all deleted}
\end{array}
\]
All clauses deleted $\Rightarrow$ \textbf{satisfiable} ($P\mapsto1,Q\mapsto0,R\mapsto0$), so the original formula is \textbf{not valid}. (If a case split is needed, show both $L$/$\neg L$ columns; reaching $\square$ in \emph{every} branch $=$ unsatisfiable $=$ theorem.)
\end{plain}

\section{Resolution --- 44\% (the headline algorithm)}

\begin{must}[title=Resolution refutation]
To prove $A$: (1) negate $A$; (2) convert $\neg A$ to clauses; (3) repeatedly apply the \textbf{resolution rule} to clauses with complementary literals; (4) derive $\square$ $\Rightarrow$ $\neg A$ unsatisfiable $\Rightarrow$ $A$ is a theorem.
\[
\frac{B\vee A \qquad \neg B\vee C}{A\vee C}
\qquad\text{set form}\qquad
\frac{\{B,A_1,\dots,A_m\}\quad \{\neg B,C_1,\dots,C_n\}}{\{A_1,\dots,A_m,C_1,\dots,C_n\}}
\]
\end{must}

\begin{plain}{Worked (propositional): prove $P\wedge Q\imp Q\wedge P$}
Negated clauses: $\{P\}$, $\{Q\}$, $\{\neg Q,\neg P\}$.
\[
\frac{\{P\}\quad\{\neg P,\neg Q\}}{\{\neg Q\}}\qquad
\frac{\{Q\}\quad\{\neg Q\}}{\square}
\]
$\square$ reached $\Rightarrow$ theorem. \textbf{Tip:} suppress repeated literals ($\{Q,Q\}=\{Q\}$); resolving a clause with its own tautology leads nowhere.
\end{plain}

\subsection{Unification (needed for first-order resolution)}

\begin{must}[title=Unification algorithm (on term trees)]
Find the \textbf{most general unifier} (MGU) $\theta$ with $t_1\theta=t_2\theta$. Cases:
\begin{itemize}[leftmargin=*,itemsep=0pt]
  \item \textbf{variable $x$ vs term $t$:} if $x=t$, return $[\,]$; \textbf{occurs check} --- if $x$ occurs in $t$, \emph{fail}; else return $[t/x]$.
  \item \textbf{constant vs constant:} equal $\Rightarrow[\,]$; differ $\Rightarrow$ fail.
  \item \textbf{pair $(t_1,t_2)$ vs $(u_1,u_2)$:} unify $t_1,u_1$ as $\theta_1$; unify $t_2\theta_1,u_2\theta_1$ as $\theta_2$; return $\theta_1\circ\theta_2$. Any failure $\Rightarrow$ not unifiable.
\end{itemize}
\end{must}

\begin{plain}{Worked unifications}
$f(x,b)$ with $f(a,y)$: unify $x,a\Rightarrow[a/x]$; then $b,y\Rightarrow[b/y]$. \ \textbf{MGU} $[a/x,b/y]$.\\
$f(x,x)$ with $f(a,b)$: $[a/x]$ then $a$ vs $b$ \textbf{fail}.\\
$f(x,x)$ with $f(y,g(y))$: $[y/x]$ then $y$ vs $g(y)$ --- \textbf{occurs check fails}.
\end{plain}

\begin{gotcha}
Skip the \textbf{occurs check} and you can ``prove'' non-theorems (it builds infinite terms / unsound proofs --- this is why Prolog, which omits it, is not a sound theorem prover). \textbf{Rename variables apart} between the two clauses before resolving --- but renaming is \emph{not} part of unification itself.
\end{gotcha}

\subsection{First-order resolution \& factoring}

\begin{must}
\textbf{Binary resolution} (with MGU $\sigma$ of $B$ and $D$):
\[
\frac{\{B,A_1,\dots,A_m\}\quad\{\neg D,C_1,\dots,C_n\}}{\{A_1,\dots,A_m,C_1,\dots,C_n\}\sigma}\quad(B\sigma=D\sigma)
\]
\textbf{Factoring} (needed for completeness --- one resolution step removes only \emph{one} complementary pair):
\[
\frac{\{B_1,\dots,B_k,A_1,\dots,A_m\}}{\{B_1,A_1,\dots,A_m\}\sigma}\quad(B_1\sigma=\dots=B_k\sigma)
\]
\end{must}

\begin{useful}
\textbf{Strategies:} \emph{set of support} --- every step must involve (a descendant of) the negated conclusion. \emph{Linear resolution} --- each resolvent becomes a parent of the next step; this is the shape of \textbf{Prolog} execution (linear resolution on \emph{Horn} clauses, left-to-right, depth-first). Delete \emph{subsumed} and tautologous clauses to stop the prover flooding.
\end{useful}

\section{Binary Decision Diagrams --- 38\%}

\begin{must}[title=ROBDD = canonical form]
A BDD is an if-then-else DAG over the propositional variables. Write $x_P Y$ for ``\textbf{if $P$ then $X$ else $Y$}'' (Shannon expansion). Three conditions make it canonical:
\begin{itemize}[leftmargin=*,itemsep=0pt]
  \item \textbf{ordering} --- variables tested in a fixed order; if $P$ before $Q$ then $P<Q$ (each var at most once per path);
  \item \textbf{uniqueness} --- identical subgraphs stored once (hash by variable $+$ child pointers);
  \item \textbf{irredundancy} --- delete a node whose true- and false-children are identical.
\end{itemize}
Then: $\textbf{equivalent} \Leftrightarrow$ same BDD (compare pointers); \textbf{tautology} $\Leftrightarrow$ BDD $=\mathbf 1$; \textbf{unsatisfiable} $\Leftrightarrow$ BDD $=\mathbf 0$.
\end{must}

\begin{must}[title=Building the BDD of a connective (the ``apply'' recursion)]
To combine $Z={}_X P\,_Y$ and $Z'={}_{X'}P'\,_{Y'}$ under a connective:
\begin{enumerate}[leftmargin=*,itemsep=0pt]
  \item recursively build BDDs for the sub-formulas, then combine top variables;
  \item \textbf{$P=P'$:} new node on $P$ with children $\mathrm{apply}(X,X')$ and $\mathrm{apply}(Y,Y')$;
  \item \textbf{$P<P'$:} branch on $P$, pairing each child of $P$ with the \emph{whole} $Z'$ (treat $Z'$ as $_{Z'}P\,_{Z'}$); symmetric if $P>P'$;
  \item \textbf{simplify} (uniqueness + irredundancy) at every step; if both children equal, return the child.
\end{enumerate}
Use a \textbf{hash table} of (op, input BDDs) $\to$ result so each combination is computed once.
\end{must}

\begin{gotcha}
\textbf{Variable ordering is everything}: a good order gives a small BDD, a bad order can be exponentially larger. State the ordering you use and test variables in it consistently. Negation is cheap: copy the BDD and swap the $\mathbf 0/\mathbf 1$ leaves.
\end{gotcha}

\section{Decision procedures: Fourier--Motzkin \& SMT --- 12\%+}

\begin{must}[title=Fourier--Motzkin variable elimination (linear arithmetic over $\mathbb{R}/\mathbb{Q}$)]
Decide satisfiability of a conjunction of linear inequalities. To eliminate variable $x$:
\begin{enumerate}[leftmargin=*,itemsep=0pt]
  \item rewrite each constraint mentioning $x$ as a \textbf{lower bound} ($x\ge\ell$) or an \textbf{upper bound} ($x\le u$);
  \item add, for \emph{every} (lower, upper) pair, the combined constraint $\ell\le u$ (with $x$ gone);
  \item keep constraints not mentioning $x$; \textbf{repeat} for the next variable.
\end{enumerate}
End with a contradiction such as $0\le c$ with $c<0$ $\Rightarrow$ \textbf{unsatisfiable}; otherwise satisfiable. (This is \emph{quantifier elimination} for $\exists$.)
\end{must}

\begin{plain}{Worked: Fourier--Motzkin on a 4-constraint system}
Constraints: $x\le y,\ x\le z,\ -x+y+2z\le 0,\ -z\le -1$.\\
Third gives lower bound $x\ge y+2z$; uppers are $x\le y,\ x\le z$. Pair them:
$y+2z\le y\Rightarrow z\le 0$;\ \ $y+2z\le z\Rightarrow z\le -y$. With $-z\le-1$ (i.e.\ $z\ge1$) we have $z\ge1$ and $z\le0$: combine $\Rightarrow 0\le -1$. \textbf{Contradiction $\Rightarrow$ unsatisfiable.}
\end{plain}

\begin{useful}
\textbf{SMT $=$ DPLL $+$ theory solver.} Treat each atomic arithmetic fact (e.g.\ $a<0$) as a propositional letter; DPLL handles the Boolean structure, and a decision procedure checks whether the chosen literals are consistent in the theory. If not, it feeds back a \emph{blocking clause} ($\neg(a<0)\vee\neg(a>0)$) and DPLL backtracks --- \textbf{over-approximation $+$ counterexample-driven refinement}. (Z3 is the canonical SMT solver.) \emph{Note:} pure-literal deletion is unsafe in SMT --- literals are connected through the theory.
\end{useful}

\section{Sequent calculus, modal logic \& tableaux --- 19\%/12\%}

\begin{must}[title=How to do a sequent proof (work backwards)]
Goal sequent $\Gamma\seq\Delta$ means ``if all of $\Gamma$ hold then some of $\Delta$ holds.'' Start from the formula you want as $\seq A$ and apply rules \textbf{upwards}, breaking connectives down, until every leaf is a \textbf{basic sequent} $A,\Gamma\seq A,\Delta$ (same formula both sides). \emph{Left} rules act on assumptions, \emph{right} rules on goals (full set on the last page).
\end{must}

\begin{gotcha}
Apply the quantifier rules with a \textbf{proviso first}: $\forall r$ and $\exists l$ require ``$x$ not free in the conclusion,'' so use them \emph{before} $\forall l$/$\exists r$ introduce terms that could capture $x$. $\forall l$ and $\exists r$ may be needed \emph{twice} (keep a copy of the quantified formula). Rename to avoid clashes.
\end{gotcha}

\begin{useful}
\textbf{Modal logic (Kripke).} A frame is $(W,R)$: worlds $+$ accessibility. $w\Vdash\square A$ iff $A$ holds in \emph{all} $v$ with $R(w,v)$; $w\Vdash\Diamond A$ iff in \emph{some}. $\Diamond A\eq\neg\square\neg A$. Systems by axiom: \textbf{K} (distribution $+$ necessitation), \textbf{T} $\square A\imp A$ (reflexive), \textbf{4} $\square A\imp\square\square A$ (transitive), \textbf{B} $A\imp\square\Diamond A$ (symmetric); \textbf{S4}$=$T$+$4, \textbf{S5}$=$T$+$4$+$B (equivalence relation). In S4, $\square\square A\eq\square A$ (operator strings collapse to length $\le$ few). \textbf{S4 sequent rules} ($\square r,\Diamond l$) \emph{discard} all non-$\square$/$\Diamond$ formulas --- so apply them \emph{first}, or you lose information and the proof fails.\\[3pt]
\textbf{Tableaux} $=$ a streamlined, refutation sequent calculus (only left rules; prove $A$ by refuting $A\seq$). \emph{Free-variable tableaux} add unification --- but then you \textbf{must Skolemise} (mixing unification with un-Skolemised $\exists$ is unsound).
\end{useful}

% ============================ ONE-PAGE RULE REFERENCE ============================
\newpage
\section*{One-page rule reference}
\begin{footnotesize}
\begin{multicols}{2}

\textbf{Propositional equivalences.}\\
$\neg\neg A\eq A$;\ \ $A\imp B\eq\neg A\vee B$;\ \ $A\leftrightarrow B\eq(A\imp B)\wedge(B\imp A)$.\\
de Morgan: $\neg(A\wedge B)\eq\neg A\vee\neg B$,\ $\neg(A\vee B)\eq\neg A\wedge\neg B$.\\
distrib: $A\vee(B\wedge C)\eq(A\vee B)\wedge(A\vee C)$ and dual.\\
simplify: $A\wedge\T\eq A$, $A\vee\F\eq A$, $A\wedge\neg A\eq\F$, $A\vee\neg A\eq\T$.

\textbf{Clause form.} negate $\to$ elim $\leftrightarrow,\imp$ $\to$ NNF $\to$ Skolemise ($\exists y$ under $\forall\bar x$ $\mapsto f(\bar x)$) $\to$ drop $\forall$ $\to$ CNF $\to$ clauses. $\square=\{\}=\F$.

\textbf{DPLL.} del tautologies; unit-propagate $\{L\}$ (drop clauses with $L$, drop $\neg L$); del pure-literal clauses; $\square\Rightarrow$ unsat; empty set $\Rightarrow$ sat; else case-split on $L$.

\textbf{Resolution.} \ $\dfrac{\{B,\bar A\}\ \ \{\neg D,\bar C\}}{\{\bar A,\bar C\}\sigma}\ (B\sigma{=}D\sigma)$.\\
\textbf{Factoring} \ $\dfrac{\{B_1,\dots,B_k,\bar A\}}{\{B_1,\bar A\}\sigma}\ (B_1\sigma{=}\dots{=}B_k\sigma)$.

\textbf{Unification (MGU).} var $x$ vs $t$: occurs-check then $[t/x]$; const vs const: equal or fail; pair: unify firsts $\theta_1$, then $t_2\theta_1$ vs $u_2\theta_1$, compose.

\textbf{BDD.} $x_P Y\equiv$ if $P$ then $X$ else $Y$. Canonical when ordered + unique + irredundant. apply: $P{=}P'$ branch on $P$ pairing children; $P{<}P'$ pair $P$'s children with all of $Z'$. Equiv $\Leftrightarrow$ same DAG; valid $\Leftrightarrow\mathbf1$; unsat $\Leftrightarrow\mathbf0$.

\textbf{Fourier--Motzkin.} per var: split into $x\ge\ell$ / $x\le u$; add $\ell\le u$ for every pair; drop $x$; repeat; $0\le c<0\Rightarrow$ unsat.

\columnbreak

\textbf{Sequent rules} ($\Gamma\seq\Delta$; basic $A,\Gamma\seq A,\Delta$).
\[\dfrac{\Gamma\seq\Delta,A}{\neg A,\Gamma\seq\Delta}\rn{\neg l}\quad
\dfrac{A,\Gamma\seq\Delta}{\Gamma\seq\Delta,\neg A}\rn{\neg r}\]
\[\dfrac{A,B,\Gamma\seq\Delta}{A\wedge B,\Gamma\seq\Delta}\rn{\wedge l}\quad
\dfrac{\Gamma\seq\Delta,A\ \ \Gamma\seq\Delta,B}{\Gamma\seq\Delta,A\wedge B}\rn{\wedge r}\]
\[\dfrac{A,\Gamma\seq\Delta\ \ B,\Gamma\seq\Delta}{A\vee B,\Gamma\seq\Delta}\rn{\vee l}\quad
\dfrac{\Gamma\seq\Delta,A,B}{\Gamma\seq\Delta,A\vee B}\rn{\vee r}\]
\[\dfrac{\Gamma\seq\Delta,A\ \ B,\Gamma\seq\Delta}{A\imp B,\Gamma\seq\Delta}\rn{\imp l}\quad
\dfrac{A,\Gamma\seq\Delta,B}{\Gamma\seq\Delta,A\imp B}\rn{\imp r}\]
\[\dfrac{A[t/x],\Gamma\seq\Delta}{\forall x\,A,\Gamma\seq\Delta}\rn{\forall l}\quad
\dfrac{\Gamma\seq\Delta,A}{\Gamma\seq\Delta,\forall x\,A}\rn{\forall r}^{*}\]
\[\dfrac{A,\Gamma\seq\Delta}{\exists x\,A,\Gamma\seq\Delta}\rn{\exists l}^{*}\quad
\dfrac{\Gamma\seq\Delta,A[t/x]}{\Gamma\seq\Delta,\exists x\,A}\rn{\exists r}\]
$^{*}$\,$x$ not free in the conclusion. \ Weakening: $\dfrac{\Gamma\seq\Delta}{A,\Gamma\seq\Delta}$, $\dfrac{\Gamma\seq\Delta}{\Gamma\seq\Delta,A}$. \ Cut: $\dfrac{\Gamma\seq\Delta,A\ \ A,\Gamma\seq\Delta}{\Gamma\seq\Delta}$.

\textbf{Modal S4 sequent rules} ($\Gamma^{*}{=}\{\square B\in\Gamma\}$, $\Delta^{*}{=}\{\Diamond B\in\Delta\}$):
\[\dfrac{A,\Gamma\seq\Delta}{\square A,\Gamma\seq\Delta}\rn{\square l}\quad
\dfrac{\Gamma^{*}\seq\Delta^{*},A}{\Gamma\seq\Delta,\square A}\rn{\square r}\]
\[\dfrac{A,\Gamma^{*}\seq\Delta^{*}}{\Diamond A,\Gamma\seq\Delta}\rn{\Diamond l}\quad
\dfrac{\Gamma\seq\Delta,A}{\Gamma\seq\Delta,\Diamond A}\rn{\Diamond r}\]
($\square r,\Diamond l$ discard non-$\square$/$\Diamond$ formulas --- apply first.)

\textbf{Modal axioms.} K: $\square(A\imp B)\imp(\square A\imp\square B)$ $+$ necessitation $\frac{A}{\square A}$; \ T $\square A\imp A$ (refl); \ 4 $\square A\imp\square\square A$ (trans); \ B $A\imp\square\Diamond A$ (symm). S4$=$T4, S5$=$T4B. $\Diamond A\eq\neg\square\neg A$.

\textbf{FOL quantifier equivalences.} $\neg\forall x A\eq\exists x\neg A$; $\neg\exists xA\eq\forall x\neg A$; pull out if $x\notin B$: $(\forall xA)\wedge B\eq\forall x(A\wedge B)$, etc.

\end{multicols}
\end{footnotesize}

\end{document}
